%!TeX root=../carlo_master.tex
%("dica" para o editor: este arquivo e parte de um documento maior)

% Palavras-chave / keywords sao obrigatorias em pt-BR e en, antes do resumo.
\palavraschave{Classificação, Teoria da decisão, Aprendizado sensível a custo,
               Otimização de limiar, Calibração de probabilidade,
               Teoria do aprendizado estatístico,
               Minimização de risco empírico}

\keywords{Classification, Decision theory, Cost-sensitive learning,
          Threshold optimization, Probability calibration,
          Statistical learning theory, Empirical risk minimization}

% Resumo (pt-BR) obrigatorio: paragrafo unico, no maximo 500 palavras,
% apresentando objetivos, metodos, resultados e conclusoes.
\resumo{
A classificação supervisionada pode ser decomposta em uma etapa de inferência, que estima probabilidades condicionais de classe, e uma etapa de decisão, que mapeia essas probabilidades em ações. Sob matrizes de custo assimétricas ou métricas de desempenho não decomponíveis, como $F_\beta$, a regra padrão $\arg\max$ é, em geral, subótima. Além disso, a regra analítica Bayes-ótima, quando existe, requer probabilidades calibradas, que classificadores reais não produzem de forma consistente.

Esta dissertação trata a regra de decisão como um componente aprendível, ajustado em uma partição de validação separada, em vez de ser obtido diretamente da fórmula analítica plug-in. A classe de decisão é pequena (um limiar escalar no caso binário e uma partição do simplex com poucos parâmetros no caso multiclasse, com dimensão polinomial no número de classes $K$). Assim, a otimização direta é, em princípio, aprendível a partir de um orçamento modesto de rótulos. Nos experimentos, entretanto, o ganho obtido só se mostra confiável quando se realiza validação cruzada sobre todo o conjunto de treinamento.

No cenário binário, derivamos o limiar analítico de Elkan, caracterizamos sua fragilidade sob má calibração e comparamos doze estratégias em quatorze conjuntos de dados, dez classificadores-base e dez objetivos.

No cenário multiclasse, organizamos as estratégias de otimização de partições em uma taxonomia unificada, calculamos suas dimensões de Natarajan, introduzimos uma parametrização de Voronoi por protótipos cujas regras, sob uma restrição de rótulo dominante, atingem risco populacional estritamente menor que o de qualquer regra de deslocamento paralelo em uma testemunha construída com perda simétrica, junto com um amostrador que preserva a restrição, e comparamos as estratégias em onze conjuntos de dados, dez classificadores-base e doze objetivos. Uma análise de sensibilidade em $K$ sobre construções sintéticas mostra que o excesso de risco da regra analítica mal especificada, relativo ao risco do oráculo, diminui à medida que $K$ cresce, em consonância com um mecanismo candidato de diluição dimensional. Todas as estratégias estão implementadas na biblioteca Python de código aberto \texttt{threshold\_opt}.
}

\abstract{
Supervised classification decomposes into an inference stage that estimates conditional class probabilities and a decision stage that maps them to actions. Under asymmetric cost matrices or non-decomposable performance metrics such as $F_\beta$, the standard $\arg\max$ rule is generally suboptimal, and the analytical Bayes-optimal rule, when it exists, requires calibrated probabilities that real classifiers do not reliably produce.

This dissertation treats the decision rule as a learnable component, fitted on a held-out validation partition rather than read off the analytical plug-in formula. The decision class is small (a scalar threshold in binary, a low-parameter simplex partition in multiclass with dimension polynomial in the number of classes $K$), so direct optimization is in principle learnable from a modest label budget; the realized gain in the experiments, however, is reliable only under cross-validation over the full training pool.

In the binary setting we derive the analytical Elkan threshold, characterize its fragility under miscalibration, and benchmark twelve strategies on fourteen datasets, ten base classifiers, and ten objectives.

In the multiclass setting we organize the partition-optimization strategies into a unified taxonomy, compute their Natarajan dimensions, introduce a Voronoi prototype parameterization whose rules, under a dominant label constraint, attain strictly lower population risk than any parallel-shift rule on a constructed symmetric-loss witness, together with a sampler that preserves the constraint, and benchmark the strategies on eleven datasets, ten base classifiers, and twelve objectives. A K-sensitivity analysis on synthetic constructions finds that the excess risk of the misspecified analytical rule, relative to the oracle risk, decreases as $K$ grows, consistent with a candidate dimensional-dilution mechanism. All strategies are implemented in the open-source Python library \texttt{threshold\_opt}.
}
